% 第 26 章。数字与教学模型：checks/check_ch26.py；图：figsrc/fig_ch26.py。
% 史料：weimar_facts；月度物价、普通预算与两种钞票的存量分别另存快照。
\chapter{一战、赔款与魏玛恶性通胀}\label{ch:weimar}

\begin{chapterguide}
\guidequestion{纸币为什么会贬值得几天就失去一半购买力？1923 年的稳定又怎样从承诺变成现实？}
\begin{guidegoals}
\item 追踪一战的战争融资怎样留下债务、货币与战后财政难题；
\item 区分赔款的财政筹措与对外转移，比较凯恩斯和俄林的论证；
\item 解释货币需求、税收滞后与汇率怎样相互推动通胀；
\item 用稳定的时序与预算数据，分析地租马克成功的条件与代价。
\end{guidegoals}
\guideprereq{\chref{ch:goldstandard}\enspace\nameref{ch:goldstandard}；\chref{ch:bullion}\enspace\nameref{ch:bullion}}
\end{chapterguide}

1923 年 11 月 15 日，德国开始发行地租马克\cite{bundesarchiv1923}。它是地租银行发行的新纸币，后来按一地租马克折合一万亿纸马克来结算。这里的马克是德国的货币单位，与称量金银的马克、汉堡的马克银行币都不同。新纸币没有让居民忽然多出粮食，却要让他们相信，收下的钱明天还买得到东西\cite{bresciani1937}。
% 据: arch_issue, bt_stablelevel

按 Bresciani-Turroni（1937）的记述，财政部长路德（Hans Luther，1879—1962）把自己的办法比作“盖房子，从屋顶盖起”。那时交易停顿，工厂关闭，城市的食物供应受到威胁。工人要求的是一种价值稳定的钱，再也不能等待\cite{bresciani1937,luther1987}。从发工资到买东西的这段时间，钱本身就在缩水。
% 据: bt_lutherroof, bt_lutherfood, bt_wantstable, luther_life

新币发行的当天，旧币还在跌。柏林的官定行情里，一美元要用约 2.5 万亿纸马克去买；五天后，要用 4.2 万亿\cite{bresciani1937}。美元是美国的货币单位，当时仍可按固定价格兑换黄金。纸马克是旧的马克纸币，地租马克是新发行的钱，两种钱要并行一段时间。货币改革已经开始，稳定还没有完成。
% 据: bt_berlin15, bt_berlin20

\chref{ch:wicksell}留下的问题在这里变得紧迫。银行不再被兑金的承诺约束，政府又有付不完的账单。可是战争结束以后，印钞并没有立即变成失控的通胀。到了危机最深的时候，人们才争相摆脱手里的纸币。为什么政府越发钱，居民越觉得钱不够用，而一张同样不能兑换黄金的新纸币，又能让局面转过来？

\section{当时的世界}\label{sec:weimar-world}

德国的兑金窗口在战争开始时就关了。1914 年 7 月 31 日下午，帝国银行自行停止把银行券兑换成黄金。8 月 4 日，法律追认了这个决定，也准许银行把国库发行的短期票据作为发钞的资产\cite{bresciani1937}。\textbf{国库券}是政府答应在较短期限内还钱的债据。银行把它买下来，先扣掉利息、交出现钱，就叫贴现。往后政府缺钱时，可以先借短债，再把债据送进中央银行。
% 据: bt_stopgold

这条路把财政部与印钞机接在了一起。财政部给工厂、军队和运输部门付款，银行得到国库券，帝国银行用自己发行的银行券换进国库券。它的资产多了一笔政府欠款，负债多了一笔公众持有的货币。\chref{ch:money}的资产负债表仍然平衡，可是资产里不再必须有同额的黄金。
% 据: bt_stopgold

政府本想用长期公债把短债接走。战时按期发行公债，居民把存款或现钞交给政府，政府用筹到的钱偿还到期短债。可是到了战争后半期，新公债筹来的钱赶不上短债的增加。最后一笔战时公债发行以后，1918 年 9 月底，短债仍约有 480 亿马克\cite{bresciani1937}。购买公债并不会把战争消耗掉的粮食、煤炭与劳动力补回来，它只决定眼下谁少花、以后谁收钱。
% 据: bt_warshort

德国把很少的战争开支交给当期普通收入。按 Bresciani-Turroni 汇集的账，战时公债认购额约 980 亿马克；整个战争期间，普通收入只覆盖约八分之一的开支\cite{bresciani1937}。普通收入不包括借债，也不能直接等同于一种税。这个比例说明，停战时接手财政的政府，面对的不只是复员，还有战时已经写下的付款承诺。
% 据: bt_warloans

增税有制度上的困难，也有分配上的阻力。战前财政收入分在帝国、各邦与地方之间，中央缺少容易随开支扩大而增加的税源。战后，复员、救济失业者、供应食物又要花钱。税务行政运转不灵，有势力的利益集团反对征收，税收因而不足\cite{bresciani1937}。
% 据: bt_fiscalstructure, bt_postwar, bt_taxopposition

削减支出也意味着把损失交给具体的人。少发救济，失业者当下就少一份收入；提高铁路运价，使用铁路的人立刻多付钱。Bresciani-Turroni 认为，革命后的政府担心，大刀阔斧的财政改革会进一步震动已经衰弱的政治经济秩序\cite{bresciani1937}。印钞把立即可见的争执往后推，却让持币的人先承担购买力损失。
% 据: bt_hardcuts, bt_postwar

\begin{compare}{征税、借债与发钞}{征税或向公众借债}{向央行借钱并发钞}
\cmprow{眼前的付款}{税款减少私人可支配收入；公债吸收愿意暂时不花的钱}{政府先付出新货币，私人手里多了名义余额}
\cmprow{以后留下什么}{税收不留还本义务；公债留下本金与利息}{公众持有货币，央行持有政府债权}
\cmprow{谁承担代价}{纳税人或将来偿债的收入提供者}{货币贬值时，持币者与来不及调价的人}
\cmprow{限制在哪里}{征税能力、政治支持与公众认购意愿}{公众愿意持有多少货币，以及发行者是否受约束}
\cmpsum{三种办法都不能省掉真实资源的消耗。发钞改变了付款的时点，也改变了谁先承受损失。}
\end{compare}

失控并不是停兑以后的直线延伸。1920 年初到 1921 年春，德国的短债与钞票仍在增加，批发物价却没有同比例上涨\cite{bresciani1937}。\textbf{批发物价指数}衡量一组批发商品价格的共同变化，与居民的生活费用不是同一个指标。\figref{fig:weimar_prices}的平缓一段很要紧。战争融资留下了危险的条件，却不能单独解释恶性通胀为什么在后来加速。
% 据: bt_1920counter

\fig{weimar_prices}{德国的物价没有沿一条斜率上涨。(a) 看 1920—1922 年月度批发物价怎样经历平台，再加速上升。(b) 看 1923 年的陡升与 1924 年的平台，纵轴为对数，单位始终保留旧纸马克；虚线标出 1923 年 11 月。指数以 1913 年$=100$。相邻月度值的变化不能当作某一天的涨幅。数据：Sargent（1982）表 G1，转录自 Young（1925）所汇集的官方统计。}{fig:weimar_prices}

\section{赔款要过两道关}\label{sec:weimar-transfer}

停战没有替政府划定一张可以一次付清的账单。赔款既加重财政负担，也使未来要付多少、何时付变得不确定\cite{sargent1982}。凯恩斯（John Maynard Keynes，1883—1946）在 1919 年离开巴黎和会。按他自己的序言，直到 6 月 7 日，他仍是英国财政部在和会的代表；他辞职，是因为不再相信和约草案会得到实质修改\cite{keynes1919,keyneslife2026}。
% 据: k19_resign, s_reparationuncertainty, keynes_life

凯恩斯的批评不只是账单太大。他问的是，欧洲靠什么继续生产和交换。法国煤矿遭到破坏，赔偿这种损失有明确的根据；和约也确实把法国煤矿的破坏列为取得萨尔煤矿的补偿理由。可是煤、铁、运输与市场互相依赖。如果抽走德国的生产条件，又要求它持续交付，实际可付的资源就会受到限制\cite{keynes1919}。
% 据: k19_france, k19_investexport, k19_coalnetwork

一国能在国内储蓄，不等于它能把同样的数额送到国外。凯恩斯写道：“德国劳动每年能为本国的资本改良生产多少剩余，在理论上和实践上，都不是它每年能向国外缴纳多少贡赋的尺度。”\cite{keynes1919}修一幢国内使用的房子，与生产外国人愿意购买的商品，需要不同的工作、设备与销路。这是他后来讨论\textbf{转移问题}的起点，也接上了\chref{ch:periphery}留下的赔款问题。
% 据: k19_investexport

常被提起的 1320 亿金马克，并不是德国立刻交出的现金。\textbf{金马克}是按战前黄金价值计价的单位，用来防止付款义务随着纸马克贬值而缩小。1921 年的付款表列出这个总额，又规定扣除已经支付和以后应获抵免的部分。债券也分批安排，其中 A、B 两类分别为 120 亿和 380 亿金马克\cite{london1921}。其余 C 类债券，要等付款资金足以支持相应本息时才由委员会发出。债务总额、某年的应付额与实际交付额，不能放在同一栏比较。
% 据: london_total, london_a, london_b, london_c

\begin{intuition}[收到了马克，还没有付到国外]
设德国政府收齐一笔马克税款，国内财政筹措的工作已经完成。可是外国收款人要用的是能在自己那里付款的钱。德国必须交付外币、黄金或约定的实物。政府可以向出口商购买外汇，也可以借到外汇；它不能只把德国税单上的金额改成外国货币的名字，就完成对外支付。
\end{intuition}

1929 年，凯恩斯把两件事分得很清楚。\textbf{财政筹措问题}是在德国居民手里征到钱，\textbf{转移问题}是把收到的德国货币变成外币\cite{keynes1929}。这场论战发生在恶性通胀结束以后，讨论的是道威斯计划下的赔款安排。它有助于看清先前的支付困难，却不是 1923 年政策制定者已经使用的共同理论。
% 据: k29_two

把其他对外收入暂时搁开，这道关可以写成一张外币账。记$R^{*}$为当期向外支付的赔款，$K^{*}$为当期净借入的外国资金。出口收入是$P_X^{*}X$，进口支出是$P_I^{*}Q$。星号表示都用同一种外币计价，$X$与$Q$是出口、进口数量。于是
\begin{keyeq}
R^{*}=\underbrace{P_X^{*}X-P_I^{*}Q}_{\text{贸易顺差}}+K^{*}.\label{eq:weimar-transfer}
\end{keyeq}
这只是为讲清支付渠道而简化的国际收支恒等式，省略了其他服务、收益与储备变动。它不说明顺差怎样出现。净借款能暂时补上差额；没有新的借款，支付就需要增加出口收入或减少进口支出。征税只完成了国内收入的转移，贸易怎么调整，还要另作解释。

\subsubsection{要不要再降一次工资}
凯恩斯最担心的是，外国市场不会自动吸收德国多出来的货。德国人少花钱，只有其中买进口品的部分会直接改善贸易余额。原来生产国内投资品的人要转去生产出口品，还得找到买主。若靠降价扩大销量，出口数量涨得多，收入也未必涨得多。他的例子是，价格降低 10\%，销量增加 20\%，出口收入只增加 8\%\cite{keynes1929}。
% 据: k29_demand, k29_borrow

出口商要降价而仍能经营，成本也得跟着降。在凯恩斯的判断里，德国既难以突然比外国更有效率，也难以得到更便宜的资本，压力就会落到按黄金计的工资上。他把差别说成：“财政筹措问题取决于德国人民的财富与繁荣；转移问题取决于其产业在国际市场上的竞争地位。”\cite{keynes1929}纳税已经让居民少花一次钱，再压低工资，又让他们承受一层损失。
% 据: k29_wages, k29_goldwages

俄林（Bertil Ohlin，1899—1979）的反驳从收款人开始。钱离开德国，同时也到了别人手里，外国的需求因此会变。德国减支、收款国增支，会直接改变两国对贸易品的需求，也会改变对国内服务与商品的需求\cite{ohlin1929,nobelohlin}。只盯着一条不变的外国需求曲线，就漏掉了购买力已经转移这件事。
% 据: ohlin_life, o29_imports, o29_indirect

国内商品的需求变化，也能把生产推向对外部门。付款国对本地商品的需求减少，劳动与资本可以更多地用于出口或替代进口。收款国对本地商品的需求增加，资源转去满足它，出口增长就会放慢，进口需求会上升。俄林因此认为，贸易余额可以在出口价格没有大幅下降时改善。他也承认，国际贸易品与国内商品之间的相对价格可能变动\cite{ohlin1929}。
% 据: o29_indirect, o29_prices

\begin{example}{谁少花，谁多花}\label{ex:weimar-transfer}
俄林举过一国借入 100 百万马克，其中 20 百万直接增加贸易品需求、其余 80 百万增加本地商品需求的例子\cite{ohlin1929}。把方向倒过来，设德国向外国转移 100 百万金马克等值的购买力。这是教学设想。再假定德国少花的钱有 20\%原用于进口，收款国多花的钱有 20\%用于购买德国商品。（1）直接效应使德国贸易余额改善多少？（2）其余部分是否已经证明无法转移？
\solution （1）德国进口减少 20 百万，出口增加 20 百万，直接改善 40 百万金马克等值的贸易余额。（2）还差 60 百万，但它不是“消失”的购买力。德国的本地商品少了需求，外国的本地商品多了需求，资源配置与价格还会调整。这个算术没有证明最后一定能改善 60 百万，也没有证明一定不能。它把问题缩到了调整要多久、生产能否转移、进口限制会不会挡路。
\end{example}
% 据: o29_direct

两边都没有给出一项不带条件的保证。俄林明确承认，大额转移会遇到困难；如果德国出口一增加，贸易保护就加强，调整可能受阻。俄林还认为，要形成出口顺差，德国几乎肯定需要让工资和物价有所下降，尤其是面向本国市场的产业。若继续快速改进生产、提高效率，名义工资又保持不变，也许可以避开部分困难。凯恩斯也承认，对外需求存在反馈，只是不相信德国能分到很大的一部分。他在答复里又加上制度条件，德国借款减少时，道威斯计划的转移保护条款可能同时让赔款减少\cite{ohlin1929,keynes1929b}。
% 据: o29_concession, o29_wageconcession, o29_productivityoffset, k29_reply, k29_repercussion

两人还同意，有组织的实物交付能绕过部分转移困难，但实施会遇到阻力\cite{keynes1929b}。
% 据: k29_inkind

两人对外债数据的读法也不同。凯恩斯认为，德国靠借外债暂时解决转移问题，借款停止时就得重新分配生产。俄林反问，借款本来就改变了德国与贷款国的购买力，不能拿有借款时的贸易形态，直接推断无借款时的能力\cite{keynes1929,ohlin1929}。所以借到了外币不能证明转移永远容易，没有出口顺差也不能证明转移永远不可能。
% 据: k29_borrow, o29_imports

\begin{compare}{赔款怎样变成一国的负担}{凯恩斯强调的困难}{俄林强调的调整}
\cmprow{先问什么}{德国怎样找到额外出口的买主}{购买力转给外国后，两国的需求怎样变}
\cmprow{价格}{扩大出口可能要降价，并压低黄金工资}{需求曲线也移动，出口价格不一定要大幅下降}
\cmprow{生产}{本国投资品不能直接变成可出口的商品}{两国对本地商品的需求变化，会推动资源重新配置}
\cmprow{各自承认什么}{国外需求有反馈；实物交付能绕过部分困难}{国内产业的工资与物价几乎肯定要有所下降；保护主义可能阻断出口}
\cmprow{证据的限制}{有外债流入时，实际所需的降价幅度看不清}{有外债流入时，购买力转移的净效应也被混在一起}
\cmpsum{两人争的是贸易怎样调整、代价有多大。把“已经征到税”等同于“已经付得出外汇”，会把争论抹掉。}
\end{compare}

赔款因此能通过多条路压在货币上。政府需要筹款，外币支付又牵动汇率；未来负担的不确定，还会改变公众是否愿意继续持有马克。可是偿付义务不直接规定必须印钞。怎样征税、削减什么、能否借款与如何安排支付，仍然是政治和经济上的选择。
% 据: k29_two, s_reparationuncertainty

\section{一张付不清的账，怎样变成更多的钞票}\label{sec:weimar-fiscal}

1923 年的鲁尔占领，把收支两头同时推向坏处。法国占领工业区后，德国政府支持消极抵抗，也就是拒绝配合占领当局、停工并由政府补助。煤税、关税和铁路收入受到损失，工人与企业却需要继续得到钱。政府给抵抗和产业的款项，靠把国库券送到帝国银行贴现来筹措\cite{bresciani1937,sargent1982}。
% 据: bt_ruhr, s_ruhrpayments

要求银行停止贴现，就意味着财政部必须立即另找付款办法。这是当时“钱不够用”的第一种意思。物价上涨，工资、采购与补助的纸马克金额更大。政府如果照旧支出、照旧征税，又没有愿意接手债券的人，中央银行拒绝新券，账单就会拖欠。单独要求银行坚守纪律，不能回答谁来承担少付款的损失。
% 据: bt_ruhr, bt_taxopposition

把政府和中央银行的账合起来，可以看清这个限制。记$D_t$为这一期的名义财政赤字，已包含利息；$B_t$为公众持有的政府债务，$H_t$为基础货币，即银行券与银行存在央行的准备金。暂不计央行对私人部门和国外资产的净购入，也不计估值变化。合并后的政府要用新增公众债务或新增基础货币填补赤字。
\begin{keyeq}
\frac{D_t}{P_t}=\frac{\Delta B_t}{P_t}+\frac{\Delta H_t}{P_t},\qquad
\frac{\Delta H_t}{P_t}=h_t-\frac{h_{t-1}}{1+\pi_t}
=\Delta h_t+\frac{\pi_t}{1+\pi_t}h_{t-1}.\label{eq:weimar-budget}
\end{keyeq}
这里$P_t$是物价水平，$h_t=H_t/P_t$是实际基础货币余额，$\pi_t=P_t/P_{t-1}-1$是当期通胀率。$\Delta$表示本期减去上期。第二个等式只把同一笔新增货币换一种写法，不能把两项当作额外的两笔收入。

第一项$\Delta h_t$，是公众愿意多持有的实际货币。第二项，是旧有货币在本期损失的购买力。用新货币取得资源，叫\textbf{铸币税}；这里已经没有铸币，名称仍然沿用。旧币购买力被物价上涨侵蚀的部分，叫\textbf{通胀税}。只有实际余额不变时，两者在这个简化账户里才相等。公众正在逃离货币时，政府即使把面额增加得很快，也未必能买到同样多的东西。

\begin{example}{多印了一倍，实际只收到多少}\label{ex:weimar-seigniorage}
设上期基础货币与物价都为 100，本期政府把基础货币增至 200，物价却升至 400。这里物价是同一篮子商品的货币价格。这些都是教学输入，忽略私人信贷等资产变化。（1）新增货币能购买多少篮商品？（2）实际余额减少多少，通胀税又是多少？
\solution （1）$\Delta H/P=(200-100)/400=0.25$篮。印钞的名义金额翻倍，实际新增购买力只有四分之一篮。（2）实际余额从$100/100=1$降为$200/400=0.5$，变化为$-0.5$。物价涨至四倍，$\pi=3$，通胀税为$3/(1+3)\times1=0.75$篮。两项相加$-0.5+0.75=0.25$，等于铸币税。旧币损失的购买力是 0.75 篮，实际余额却同时减少了 0.5 篮，政府取得的资源只有两者的差额。
\end{example}

税收还会在路上缩水。政府先按一个名义金额核定税款，隔一段时间才收到钱。物价在这段时间上涨，国库收到的真实资源就少了。拖延纳税因而变得有利；企业代扣的税款，也可能暂时留在自己手里。Bresciani-Turroni 记下，1923 年夏季，工资税平均在雇主手里停留约两个星期，才交进国库\cite{bresciani1937}。
% 据: bt_wagetax

\begin{example}{税款在途中失去多少}\label{ex:weimar-taxlag}
以工资税约两周的滞后为背景，假设物价每七天翻一番。这是教学假设，不是德国某月的实测日率。核定税款时，这笔税可买 100 篮商品。两周后收进国库，税款还能买多少？如果核定额仍不变，政府要维持原来的支出，缺口有多大？
\solution 两周经过两次翻倍，物价变为四倍。税款只能买$100/4=25$篮，缺口是 75 篮。把税率提高一点，可能很快又被收款滞后吞掉。按稳定价值核定税款、缩短收款间隔，能减少这一损失；若支出与货币发行照旧加速，调整本身也会追不上物价。
\end{example}

这条反馈在 1922 年后尤其明显。此前财政恶化推动货币贬值，此后贬值又削弱税收的真实价值，铁路收费也赶不上成本。赤字因此进一步扩大\cite{bresciani1937}。把税额改按黄金价值计算，有助于堵住漏洞；可是在最深的危机里，换算系数刚公布又已经过时。早期财政失衡推动通胀，不等于后来的赤字都可以当作独立的起因。
% 据: bt_fiscalfeedback, bt_valorize

\section{为什么人人都想把钱花掉}\label{sec:weimar-demand}

钱不够用还有第二种意思，居民拿到手的现金买不到原来的东西。一个家庭并不需要先学会货币理论，才知道等待有代价。工资刚到手，早一天买到食物，就少承受一天的贬值。商人也要补货，卖货收到的纸币若跌得比补货价格还快，账面上的盈利会变成买不回原料的亏损。
% 据: bt_lutherfood, bt_payments, bt_noconsumption

人们于是尽量缩短持有纸币的时间。能换外币的换外币，能提前买货的买货，企业逐渐改用外币或金马克报价格。Bresciani-Turroni 记述，这种计价在 1922 年夏季已扩展到许多重要行业\cite{bresciani1937}。同样的纸币在更多交易中迅速转手，银行券的面额尽管增加，愿意留下它的人却减少了。
% 据: bt_payments, bt_substitution

把钱花出去，个人能减少自己的暴露，全体却不能靠互相传递把旧币消灭。纸币到了收款人手里，收款人也想赶快花掉。价格与汇率因此要变到有人愿意暂时留下这笔钱为止。能否换成别的支付工具、是否必须用马克纳税，以及政府以后会怎样发钱，都会影响这个选择。

\textbf{货币需求}是人们愿意保留的购买力，不是他们希望收入越多越好。记$M_t$为一项口径固定的货币存量，$Y_t$为当期实际交易所对应的产出，$k_t$为愿意用货币留住的收入份额。用最简化的数量关系来写，就是
\begin{keyeq}
\frac{M_t}{P_t}=k_tY_t,\qquad
P_t=\frac{M_t}{k_tY_t},\qquad V_t=\frac1{k_t}.\label{eq:weimar-demand}
\end{keyeq}
$V_t$是收入流通速度。这个式子接着\chref{ch:bullion}的交易方程，但不能把$k_t$固定不动。物价涨得越快，留下货币的代价越高，人们愿意留住的份额往往越小。正式的动态货币需求模型留给\chref{ch:hyperinflation}。这里要学会的是，给定产出时，货币增加与需求减少都能推高物价。

汇率让这个过程更快进入日常生活。记$e$为一美元折合的纸马克数，$e$上升就是马克贬值。进口商要付美元，纸马克贬值立即推高补货成本；预料以后会更贵的人，提前买外汇。政府也要购买外汇或支付随成本上涨的账单。帝国银行再多发钱，反过来又使马克承受贬值压力\cite{bresciani1937}。
% 据: bt_pricesfeedback

\begin{compare}{哪一个先动}{先看到汇率与物价上涨}{先看到财政融资与货币增加}
\cmprow{最强的观察}{外币行情可以立刻跳动，税款与工资随后追赶}{战时开始的短债贴现，持续替政府提供付款手段}
\cmprow{能说明什么}{预期、进口成本与持币意愿能推动短期加速}{持续赤字怎样得到货币融资，发行为什么缺少硬限制}
\cmprow{不能单独说明什么}{银行为什么持续满足更大的名义付款需求}{同样的货币增幅，为什么有时伴随更快的物价上涨}
\cmprow{两条路相遇时}{贬值削弱税收、扩大名义缺口}{再融资增加货币，进一步削弱持币意愿}
\cmpsum{时间上领先不等于唯一的原因。财政安排与对货币的预期会互相作用。}
\end{compare}

当时的一种解释恰好把顺序倒过来。黑尔弗里希（Karl Helfferich，1872—1924）认为，赔款和法国的强制政策先压低马克。进口品、国内物价与工资跟着涨，公众与政府因此需要更多现钱，帝国银行才被动增加发行。他的解释，经 Bresciani-Turroni 转引\cite{bresciani1937,helfferich1969}。它最有力的事实，是物价与汇率在一些时期确实跑在钞票数量前面。
% 据: bt_helfferich, helfferich_life, s_ratio1, s_ratio2

财政解释的反问是，为什么把“缺钱”当成必须无限贴现的理由。兑金义务已经停止，中央银行不必用有限的黄金来制约新增贷款。反复满足支出需要，会使人相信下一轮还会增发，今天的货币需求便进一步下降。可是外部冲击确实损害了税源和生产，不能靠账上的货币项目把它们排除\cite{bresciani1937,sargent1982}。
% 据: bt_stopgold, bt_ruhr, bt_fiscalfeedback

一张涨在前面的汇率图，不能替任何一方结束争论。未来发钞的预期可以先影响今天的汇率，已经发生的贬值也会迫使财政和企业重新计算付款。1923 年后期，几条曲线越来越紧地相互追赶\cite{bresciani1937}。能确认的是反馈在加强；仅凭月度先后，分不出每条路各贡献了多少。
% 据: bt_pricesfeedback

\section{数据检验：钱多了，购买力却少了}\label{sec:weimar-test}

回到路德面对的缺钱问题，数字里确实出现了这个反差。从 1922 年 1 月到 1923 年 7 月，批发物价约变为原来的两千倍，帝国银行券约变为原来的 380 倍\cite{sargent1982}。用这个价格指数折算，银行券所代表的实际购买力只剩起点的不到五分之一。到 8 月，两个倍数都更大，真实余额仍然很低。
% 据: s_ratio1, s_ratio2

这能说明公众在节省使用帝国银行券，不能直接量出全部货币需求。统计只包括帝国银行的银行券，没有包括银行存款、外币和形形色色的应急货币。价格用的是批发商品，不是每个人的消费篮子。产出与交易也在变化。因此$M/P$的下降只是与逃离纸币相符的证据，不能把全部变化都归给式~\eqref{eq:weimar-demand}里的$k$。
% 据: s_ratio1, s_ratio2, bt_substitution

\figref{fig:weimar_balances}把两个容易混淆的时期并排放在一起。左边是危机中银行券的数量追不上物价，右边是稳定以后两种正式纸币合计继续增加。左边的钱多而购买力少，右边的钱增加而汇率大体稳定。同一个“货币数量增加”，背后的持有意愿和资产安排可以很不一样。
% 据: s_ratio1, s_ratio2, bt_oldcirculation, bt_newcirculation

\fig{weimar_balances}{先看数量追不上物价，再看稳定后正式货币重新增加。(a) 1922 年 1 月$=1$，1923 年 7 月与 8 月的批发物价、帝国银行券倍数；纵轴为对数。两数之比是以批发物价折算的银行券实际余额相对值。(b) 三个年末或月末，两家银行的钞票分别折成百万地租马克等值后相加。旧纸马克除以$10^{12}$；不包括存款、应急货币和外币。数据：Sargent（1982）第 82 页；Bresciani-Turroni（1937）第 337 页。左图的概括倍数与右图的存量口径不同，不能接成一条总货币序列。}{fig:weimar_balances}

“每天翻倍”也要经过口径检查。逐日的市场报价可以在个别日子出现极端变化，它既不等于全国所有商品的价格，也不等于整个月的速度。1923 年 9 月到 10 月，月度批发物价指数约变为原来的 296 倍；若用三十天匀速增长作换算，大约每 3.7 天翻一番。这个换算已经极端，却仍然不同于每天翻倍。

\begin{example}{把月度指数换成能想象的速度}\label{ex:weimar-double}
表 G1 给出的 1923 年 9 月与 10 月批发物价指数，分别为$2{,}394{,}889{,}300$与$709{,}480{,}000{,}000$，基期为 1913 年$=100$\cite{sargent1982}。取三十天匀速增长作为教学换算。（1）指数之比、百分比涨幅是多少？（2）相当于多少天翻一番？（3）如果真是每天翻倍，三十天后应变为多少倍？
\solution （1）比值约 296.25，涨幅约 29,525\%。倍数要减去 1，再乘 100，才是百分比涨幅。（2）令$2^{30/d}=296.25$，得$d=30\ln2/\ln296.25\approx3.65$天。（3）$2^{30}=1{,}073{,}741{,}824$倍，远比 296 倍大。这不是实测的每日通胀路径，也不能据此认定全时期最快的一天或一个月。
\end{example}

官价与市场价也不是一回事。11 月 20 日，柏林官定美元价是 4.2 万亿纸马克；Bresciani-Turroni 转引的科隆市场报价却约为 11.7 万亿\cite{bresciani1937}。宣布一个价格以后，还要使交易者愿意照着成交。新币发行日、官价固定日与市场接受这个价格的过程，不能压成同一个瞬间。
% 据: bt_berlin20, bt_cologne20, bt_marketgap

恶性通胀对不同人的损害，并不按同一比例分摊。能迅速重报价格、持有实物或外币的人，有机会保护自己。按固定金额领取工资、养老金和利息的人，很难同样快地调整。Bresciani-Turroni 所整理的公务员工资，扣除生活费用以后，在整个通胀期低于战前水平；中高层的跌幅还比低层大\cite{bresciani1937}。名义工资上涨，不能据此说生活改善。
% 据: bt_income, bt_payments

债权人还会遭遇另一种损失。纸币债权承诺的是一个名义金额，价格上涨以后，这个金额买到的东西减少。债务人则可能以已经贬值的钱偿还。费雪的实际利率算法在这里仍然适用（式~\eqref{eq:prices-fisher}）。如果物价涨幅远远高于合同利率，本息收回以后，实际购买力可能几乎消失。这是旧合同的重新分配，不是物价稳定以后就会自动倒转的过程。

\begin{pitfall}[纸币发得太少，才造成恶性通胀]
银行券的实际购买力很少，与名义发行过多并不矛盾。居民不愿持有，实际余额就能在大量增发中收缩。继续照着已经上涨的价格补发，可能把持币意愿进一步推低。反过来，只看钞票面额也会误判稳定后的增加。要先说明统计是哪家机构的负债、价格用什么指数、有没有其他支付工具，再判断新增货币满足了什么需求。
\end{pitfall}

\section{地租马克怎样站住}\label{sec:weimar-stabilization}

路德没有等到赔款争端全部解决才开始改革。1923 年秋，政府先结束支持鲁尔消极抵抗的政策，再整顿财政、建立新发行机构。10 月 15 日的地租银行法令，由总理施特雷泽曼（Gustav Stresemann，1878—1929）与路德签署\cite{renten1923,bundesarchiv1923b}。政治上先承认旧办法难以为继，才可能让限额有意义。
% 据: chron_endpassive, law_signed, stresemann_life

地租马克的名字来自地产和企业承担的担保。法令把农业地产上的抵押与工商业的义务，变成新机构的资产。法令还要求公共金库收下新钞作为付款。可是居民不能拿一张地租马克去分走一块土地，也不能随时换到金币。它仍是纸币。法令允许以规定数额的新钞换取按金马克计价的有息债券，这和现金兑金不同\cite{renten1923,bresciani1937}。
% 据: law_collateral, law_bond, law_taxes, bt_convert

真正要观察的是，这套担保会不会伴随无限的新债。法令规定的资本与基本储备合计 32 亿地租马克，纸币不得超过这一金额。给政府的信贷限额为 12 亿，给私人经济的信贷也另有限额\cite{renten1923}。所以总发行上限、政府借款上限与当时实际发出多少，是三种不同的数。资产有一个响亮的名称，不足以替代这些限制。
% 据: law_capital, law_limit, law_credit, law_private

旧国库券的入口也被关上。法令第十九条规定：“地租银行一开始发行银行券，帝国银行就不得再贴现国库券。”\cite{renten1923}它仍允许在清偿旧券期间办理展期。条文针对的是增加政府融资的那条通道，不能读成旧钞全部注销，也不能读成所有商业信贷同时停止。
% 据: law_stop, law_rollover

新办法并非要求政府一分钱都不能再借。新机构给了一个有限的过渡额度，其中一部分用来清偿旧国库券，其余帮助政府渡过整顿期\cite{renten1923,bresciani1937}。这种安排把“今天尚有缺口”与“以后每个缺口都由央行满足”分开。12 月政府试图突破已定限制，地租银行管理层拒绝了。限额因此不只是一段公布过的文字\cite{bresciani1937}。
% 据: law_debt, law_credit, bt_decembertest

财政措施让这份拒绝有了支撑。税额和公共服务收费逐步改按稳定价值核定，政府削减支出并裁减人员。改革有一个落实过程，不能以发行日为界，假定前一天财政全面失控，后一天所有账都平了。\tabref{tab:weimar-stages}列出的，是几项互相配合的改变。
% 据: bt_valorize, bt_finaltax, s_cost

\begin{table}[!tb]\centering\small
\caption{稳定的步骤与尚未完成的工作}\label{tab:weimar-stages}
\begin{tabularx}{\linewidth}{@{}p{0.17\linewidth}YY@{}}
\toprule
时点 & 已发生的变化 & 不能据此断言什么 \\\midrule
1923 年 9 月 & 结束消极抵抗政策，减少继续支持它的支出压力 & 鲁尔经济与政府全部收入已经恢复 \\
10 月 15 日 & 颁布地租银行法令，规定发行与政府借款限制 & 新钞已经进入每个家庭，旧钞已经稳定 \\
11 月 15 日 & 新币开始发行，停止新增国库券贴现的规则付诸实施 & 所有纸币发行停止，财政当月立即完全平衡 \\
11 月 20 日以后 & 柏林美元官价固定在 4.2 万亿纸马克 & 各地自由市场立即与官价一致 \\
12 月至次年初 & 税收与支出整顿继续，普通预算明显改善 & 每个月、每个政府账户都不再赤字 \\
1924 年 & 道威斯计划重整赔款安排，货币制度进一步重建 & 外债安排是 1923 年 11 月稳定的先决条件 \\
\bottomrule
\end{tabularx}
\datasource{地租银行法令（1923）第 14—19 条；Bresciani-Turroni（1937）第 24、73、335—350 页；Sargent（1982）第 83—84 页；联邦档案馆鲁尔占领史料页。不同文本把开始发行和停止贴现记在 11 月 15 日或 16 日，见史料说明。}
\end{table}

预算表比“一夜平衡”的说法更细。以金马克计的联邦普通预算，1923 年 12 月收入约 3.3 亿，支出约 6.7 亿，仍有约 3.3 亿的缺口。次年 1 月出现盈余，2 月又有小额赤字\cite{sargent1982}。这只是联邦普通账户，不能拿来代替整个公共部门。它说明稳定与财政改善相伴，却没有说明发行当天一切都已完成。
% 据: s_budgetscope

\fig{weimar_budget}{预算在稳定后继续调整。柱子是德国联邦普通预算的收入减支出，正值为盈余、负值为赤字，单位为百万金马克；1923 年 12 月至 1924 年 6 月。它不包括全体政府与所有特别账户，11 月因支出数据空缺不画。物价与汇率的转折先于表中连续出现盈余。数据：Sargent（1982）表 G5，转引 Young 所录《经济与统计》的官方账。}{fig:weimar_budget}

稳定也没有冻结钞票存量。旧纸马克还在流通，帝国银行还为商业用途提供货币；地租银行又发出新钞。到 1924 年 3 月底，两种正式纸币合计约为 24.5 亿地租马克等值，明显超过 1923 年末\cite{bresciani1937}。公众愿意多留一些钱，正式货币又替代了一部分应急货币与外币。资产构成与货币需求同时改变，单看面额相加的增长，便容易误判。
% 据: bt_oldcirculation, bt_newcirculation, bt_noconsumption, bt_substitution

\begin{compare}{换了名字，还是换了制度}{只改面额和名称}{1923 年的稳定安排}
\cmprow{记账单位}{给旧金额删去零，方便报价}{新单位与旧单位并行，最后按固定比率结算}
\cmprow{政府融资}{仍可无限把新债交给央行}{停止新增国库券贴现，设有限过渡信贷}
\cmprow{财政收支}{原有收款滞后和支出义务不变}{税收、收费和支出同时整顿}
\cmprow{公众选择}{如果预料继续贬值，仍会迅速花掉}{愿意重新持币，正式货币逐渐替代其他支付工具}
\cmpsum{改单位能简化算术。限额能否执行、政府能否另筹收入，决定新单位的购买力能不能站住。}
\end{compare}

由这些事实还不能识别出一种政策的单独功劳。财政整顿、新币、信贷限制、汇率政策和公众预期几乎同时变化。外部支付得到喘息，也减轻了约束。可以排除的是两个过强的说法，既不能把成功只归给纸币背后的土地，也不能说必须先完全解决赔款，才可能稳住货币。道威斯计划来得更晚，初步稳定已经发生\cite{bresciani1937,sargent1982}。
% 据: bt_convert, bt_decembertest, s_dawes, law_signed

\section{稳定之后，谁还在付代价}\label{sec:weimar-cost}

稳住购买力，救不了已经损失的旧储蓄。货币停止贬值以后，居民拿到工资到花出去之间的损耗减少，交易也容易恢复。可是旧债权按名义金额缩水的损失，不会因新币出现而自动补偿。哪些债权要重估、给谁补偿，又是新的分配争执。停止通胀，解决的是继续损失的问题，不能让过去的所有合同重新兑现。
% 据: bt_lutherfood, bt_income, bt_debtrevaluation

停止容易获得的信用，也会使依赖它的企业难以周转。财政削减则直接减少一些人的收入。按 Sargent 转引的资料，铁路系统 1923 年裁减约 12 万人，1924 年又裁减约 6 万人\cite{sargent1982}。前一个数字包含稳定以前的月份，不能把全部裁员都算作 11 月改革造成。它至少提醒读者，财政表上的节约，对应着离开岗位的人。
% 据: s_cost

\begin{case}{奥地利与匈牙利 · 国联怎样帮助稳定}
奥地利的克朗也是战后贬值的纸币，不能把它与德国马克直接按面额比较。1922 年 8 月，奥地利汇率先出现转折；10 月签署的国联议定书，安排了外部贷款、财政重建和对政府借款的限制\cite{sargent1982}。一边是短期取得资源的办法，一边是今后不再靠不断发钞的约束。正式改称先令，要到 1924 年末才发生，每先令折合一万旧克朗。稳定因而早于改名字\cite{sargent1982}。
% 据: s_ausaug, s_ausloan, s_auscommit, s_ausunit

匈牙利在 1924 年的国联重建中，也把外部贷款与新中央银行的融资限制结合起来。新银行不得再向政府追加没有足额黄金或外汇担保的贷款\cite{sargent1982}。德国在危机时没有先取得同样的外部贷款，使用的是有上限的国内过渡融资。外援可以改变过渡期怎样付账，却没有免掉财政整顿的问题。
% 据: s_hungary, s_hungarybank, law_credit

奥地利的失业救济领取人数，在稳定后仍一度上升。到 1923 年 3 月约为 17 万人，此后才下降\cite{sargent1982}。战争后的产业分割、原有失业与稳定措施同时影响就业，不能从先后顺序里精确分出代价。这两个案例说明新币名称不是充分条件，也不能用来证明迅速稳定一定没有就业损失。更系统的跨国比较见\chref{ch:hyperinflation}。
% 据: s_auscost
\end{case}

德国也不能只挑一个年份判断稳定的代价。鲁尔占领与消极抵抗已经压低生产，拿 1924 年同 1923 年比较，会包含恢复生产的收益。把改革后的所有困难归给紧缩，又会漏掉此前通胀改变投资、侵蚀储蓄的后果\cite{sargent1982}。评估稳定的成本，需要说明同哪一条没有发生的路径比较；现有历史表格不能单独回答这个反事实。
% 据: s_outputcomparison

\begin{deeper}[史料说明与检验的边界]
\textbf{版本与人声。}凯恩斯 1919 年的书据古登堡电子本；1929 年的三月论文与六月论战，直接核对《经济学杂志》原期识别文本。黑尔弗里希、路德的话据 Bresciani-Turroni 的 1937 年英译修订版，未直接核他们的原书；关于政府恐惧改革冲击的判断，也是该作者对决策的解释。该版注明“ALL RIGHTS RESERVED”，本章仅作必要的短引与转述。德国国库券贴现的法律依据，另核 1923 年《帝国法律公报》的法令录文。

\textbf{两天与两种限额。}档案馆把新币发行记为 11 月 15 日，Bresciani-Turroni 在第 337 页叙述实际发行时写作 16 日，其第 23 页又把停止国库券贴现记在 15 日。正文采用 15 日作为制度启动日，不据一天之差推断政策的先后效应。法令第 2、14 条的总发行上限为 32 亿，政府信贷 12 亿；第 16 条另准向私人部门提供 12 亿信贷。Bresciani-Turroni 第 348 页注称最大发行为 24 亿，与两类信贷之和相同；不能把这个数代替法令上限。

\textbf{大数的译法。}1937 年英文本用英国旧制，milliard 为$10^9$，billion 为$10^{12}$，trillion 为$10^{18}$。Sargent 表 G2 的一个脚注又把新旧马克比率写成 billion，其他地方明确写$10^{12}$。本章一律用幂次核算。例如旧币 496 个英制 trillion，折成新单位以后是 4.96 亿，不是 496 张，也不是 496 亿。

\textbf{物价与预算。}月度批发物价来自 Sargent 表 G1，原始编制者是德国统计机关，Sargent 取自 Young 的 1925 年汇编。这里只读到转录，没有取得 Young 原表作独立复核。序列到 1924 年 6 月止，避开 7 月改用新马克的单位跳变。倍增时间是月度指数的三十天等速换算，不是日价格重建。预算表 G5 只记联邦普通收支；Sargent 正文的迅速平衡不能覆盖表中 12 月与次年 2 月仍有赤字的事实。

\textbf{单一见证与因果。}物价、预算、战后裁员、国联案例据 Sargent 所转引的统计。科隆自由汇率、两种纸币存量、公务员实际工资与路德的比喻，据 Bresciani-Turroni 一个来源。原表的覆盖范围与读数保留在快照中。两份著作各有明确的解释立场，数据与解释分开使用；政策同时变化，不据曲线共动声称已测出财政、货币需求或汇率各自的因果贡献。
% 据: bt_edition, bt_issue16, bt_stopgold, arch_issue, bt_limit24, law_capital, law_limit, law_credit, law_private, s_uniterror, bt_oldcirculation, s_budgetscope
\end{deeper}

\section{今天还剩什么}\label{sec:weimar-legacy}

\textbf{货币发行的财政背景不能省掉。}政府持续有缺口，又失去其他筹款渠道，中央银行会面对越来越强的融资压力。把它写成政府预算约束，得到的是可行性的限制；由谁调整、何时调整，还需要制度与政治的解释。更一般的财政与货币关系见\chref{ch:budget}和\chref{ch:fiscaldominance}。

\textbf{货币需求会随发行制度改变。}人们预料的钱值不值钱，会影响今天愿意留多少。1923 年的货币数量与物价没有固定比例，稳定后的货币增加也不一定意味着新一轮失控。可是用银行券除以批发物价，还只是有限口径的证据。怎样建立和检验动态需求，是\chref{ch:hyperinflation}的问题。

\textbf{对外支付需要两国一起调整。}财政部收齐本币，外币支付仍要经过贸易、借款与资产变动。凯恩斯与俄林争的不是这张账是否成立，而是贸易怎样响应购买力转移、相对价格要变多少。这个问题同样出现在外债偿还与国际金融危机中（\chref{ch:emcrises}）。

\textbf{停止通胀与重建黄金平价是两件事。}德国先让旧币停止继续下跌，再以新单位组织支付。新纸币站得住，不要求先把旧债权的购买力恢复到战前。别国却可能选择另一条路，把货币重新推回旧的黄金价值。平价怎样选、谁为调整付钱，是\chref{ch:interwar}的问题。

\begin{takeaways}
\begin{itemize}
\item 战时停兑和国库券贴现打开了货币融资通道。战后收支压力、税收能力与利益冲突让它延续，但战争不能解释后来每一段加速。
\item 赔款有财政筹措与对外转移两道关。外债可以暂时支付，购买力的跨国移动又会改变两边的需求，不能固定出口需求来推断。
\item 新增货币取得的资源，等于实际余额的增加加上旧余额被通胀侵蚀的部分（式~\eqref{eq:weimar-budget}）。公众减少持币时，名义发行更多，取得的真实资源未必更多。
\item 贬值、税收滞后、赤字融资和逃离纸币形成反馈。物价领先货币，并不能单独证明货币发行只是无害的被动反应。
\item 地租马克的成功伴随发行限额、停止新增国库券贴现和财政整顿。普通预算尚有赤字时，汇率已经转折；正式纸币后来仍在增加。
\item 稳定使继续持币不再迅速受损，却不能自动修复旧储蓄，也没有消除裁员和信用收缩的代价。
\end{itemize}
\end{takeaways}

\begin{selfcheck}
\begin{enumerate}
\item 德国政府征到本币税款，是否就解决了转移问题？假设当期没有其他对外收入，赔款为 100，贸易顺差为 40，须净借入多少外币？这些数都是教学输入。
\item 凯恩斯的出口例子中，价格下降 10\%，销量增加 20\%，出口收入为什么只增加 8\%？俄林认为这个计算漏掉了什么变化？
\item 上期基础货币与物价均为 100。本期货币为 300，物价为 500。按式~\eqref{eq:weimar-budget}求铸币税、实际余额变化与通胀税。数字为教学输入。
\item 一个居民持有固定名义金额，物价在收款以后两次翻倍。这笔钱剩下多少购买力？为什么每个人都尽快花钱，会使总体物价更难稳定？
\item 地租马克没有随时兑金的承诺，为什么仍可能稳定？为什么稳定以后正式纸币数量增加，并不立即构成反例？
\item 1923 年 12 月的普通预算仍有缺口，是否足以推翻财政整顿有助于稳定的解释？需要再看什么？
\end{enumerate}
\end{selfcheck}

\begin{answers}
\item 没有。本币要变成外币、黄金或实物交付。按式~\eqref{eq:weimar-transfer}，净借款为$100-40=60$。这没有保证借得到，也没有保证能偿还。
\item $0.9\times1.2=1.08$，收入只多 8\%。俄林强调，两国购买力变化会移动需求，也会推动资源在国内商品和贸易品之间转移，不能把需求曲线固定住。
\item 铸币税为$(300-100)/500=0.4$。实际余额由 1 降为 0.6，变化为$-0.4$；通胀率为 4，通胀税为$4/5\times1=0.8$。$-0.4+0.8=0.4$。
\item 剩四分之一。居民缩短持币时间，降低了愿意持有的实际余额；若名义存量不随之减少，就会形成对商品和外币的更多购买压力。
\item 关键是增发的约束、财政怎样付账，以及公众是否愿意持有。需求恢复和支付工具的替代，能吸收更多纸币；还要核对新增资产，不能把全部发行当作财政透支。
\item 不足以。要看缺口如何融资、以后能否减少、收支制度有何改变，也要分清普通预算与总账。有限的过渡融资，与无限的持续货币融资不同。
\end{answers}

路德所谓“从屋顶盖起”的办法，最后没有靠居民拿纸币去换土地来兑现。人们重新愿意收钱、留钱，政府也逐步把付款移出无限贴现的通道。可是 1923 年的稳定没有把旧马克推回战前的购买力。别国若坚持回到旧的黄金平价，就得让国内价格与工资去迁就它。战争后的欧洲要恢复金本位，究竟应恢复哪个平价，又由谁承担代价？
% 据: bt_lutherroof, bt_stablelevel, bt_decembertest
